\section{Question 8}
\begin{enumerate}[label=(\alph*)]
    \item
    \begin{enumerate}[label=(\roman*)]
        \item 
        By the RADIUS of CONVERGENCE Formula,
        \begin{align*}
            R &= \lim_{n\to\infty} \left| \frac{3^n + 5n}{n!} \times \frac{ (n+1)! }{ 3^{n+1} + 5(n+1) }  \right| \\
              &= \lim_{n\to\infty} \left| \frac{ \left(1 + \frac{5n}{3^n}\right)n }{ 3 + 5\left(\frac{n+1}{3^n}\right) } \right| \\
              &= \lim_{n\to\infty} \left| \frac{(1+0)n}{3 + 5(0)} \right| \to \infty \text{ as } n \to \infty.
        \end{align*}
        So the radius of curvature is infinite.
        \vvs \\
        \item 
        By the RADIUS of CONVERGENCE Formula,
        \begin{align*}
            R &= \lim_{n\to\infty} \left| \frac{ (2n^2 + e^{-n})2 }{ 2( 2(n+1)^2 + e^{-(n+1)} ) } \right| \\
            \intertext{but $e^{-n},e^{-(n+1)} \to \infty \text{ as } n \to \infty $,}
              &= \lim_{n\to\infty} \left|  \frac{ 2n^2 + 0 }{ 2(n+1)^2 + 0 } \right| \to 1.
        \end{align*}
        Thus, the radius of curvature is 1.
    \end{enumerate}
    \leavevmode \vvs \\
    \item 
    \begin{enumerate}[label=(\roman*)]
        \item 
        As $ \sinh 0 = 0 $, the condition in HB B3 4.5, p53 is satisfied and we can use the COMPOSITION RULE for Power Series.
        \vs \\
        Thus,
        \begin{align*}
           f(z) &= \sin\left( z + \frac{z^3}{3!} + \frac{z^5}{5!} + \ldots \right) \quad \text{for } z \in \C \\
             &= \left( z + \frac{z^3}{3!} + \frac{z^5}{5!} + \ldots \right) - \frac{1}{3!}\left( z + \frac{z^3}{3!} + \ldots \right)^3 + \frac{1}{5!}\left( z + \ldots \right)^5 \quad \text{for } z \in \C \\
             &= \left( z + \frac{z^3}{3!} + \frac{z^5}{5!} + \ldots \right) - \frac{1}{3!}\left( z^3 + \frac{z^5}{3!} + \ldots \right) + \frac{1}{5!}\left( z^5 + \ldots \right) \\
             &= z + \left( \frac{1}{3!} - \frac{1}{3!} \right)z^3 + \left( \frac{1}{5!} - \frac{1}{12} + \frac{1}{5!} \right)z^5 + \ldots \\
             &= z - \frac{z^5}{15} + \ldots \quad \text{for } z \in \C, \text{ as required}.
        \end{align*}
        \vvs \\
        
        {
        \color{blue}
        \item For this part, I used the series from (b)(i). Substituting $1/z$ for $z$ produces infinitely many terms in negative powers of $z$ in the Laurent expansion. So, by HB B4 2.10(c), p57, $f(1/z)$ has an essential singularity at 0.

        But I also remember looking at it and thinking I probably hadn't written enough for 3 marks!

        -- John
        \vvs


        \item 
        This is actually part (iii)
        }

        $ z^4 f(1/z) $ has a Laurent Series around 0 of
        \[ z^3 - \frac{1}{15z} + \ldots \; . \]
        Hence, by the defintion of Residues:
        \[ \Res \big(z^4 f(1/z), 0\big) = -\frac{1}{15}. \]
        Let $ \mathcal{R} $ be the simply connected region $\C$. Then $ z^4 f(1/z) $ is analytic on $ \mathcal{R} $, except for the singularity at 0 and $C$ is a simple-closed contour in $ \mathcal{R} $ not passing through 0.
        \vs \\
        Thus, by CAUCHY'S RESIDUE Theorem,
        \begin{align*}
            \int_C z^4 f\left( \frac{1}{z} \right) \dd z &= 2\pi i \Res\left( z^4 f\left( \frac{1}{z} \right), 0 \right) \\
            &= -\frac{2\pi}{15}i.
        \end{align*}
    \end{enumerate} 
    \item 
    \begin{enumerate}[label=(\roman*)]
        \item 
        Note first that $ \cosh z - \sinh z = e^{-z} $ and $ |e^{-z}| > 0 $.
        \vs \\
        If we choose
        \[ f(z) = \cosh z - e^{-z} +3 \text{ and } g(z) = \frac{\sinh z}{3}, \]
        then both $f$ and $g$ are non-constant and entire.
        \vs \\
        So,
        \begin{align*}
            |f(z) - 3g(z)| &= \left| \cosh z - e^{-z} +3 - 3\left(\frac{\sinh z}{3}\right)  \right| \\
            &= |e^{-z} - e^{-z} + 3| = 3.
        \end{align*}
        Hence, $ f(z) = \cosh z - e^{-z} +3 \text{ and } g(z) = \dfrac{\sinh z}{3} $ are both non-constant, entire, and satisfy $ |f(z) - 3g(z)| >2 $ for $z \in \C$, as required.
        \vvs \\
        \item 
        Let $ g(z) = \left( \dfrac{e^z - 1}{3} \right) $, then $g$ is entire and non-constant, with $g(0) = 0$ and 
        \[ | e^z - 3g(z) | < 2, \quad \text{for } z \in \C, \text{ as required}. \]
    \end{enumerate}
\end{enumerate}

